% !TEX root = ../main.tex

\chapter{绪论}

\section{研究背景}

随着具身智能、服务机器人、巡检机器人与救援机器人的快速发展，机器人系统已经不再满足于在规则明确、几何结构稳定的实验室环境中完成预设动作，而是需要在真实世界中面向更复杂的任务目标进行持续自主运行。无论是园区物流、楼宇巡检、家庭服务，还是室内外混合场景下的目标搜索与自主探索，机器人都必须具备稳定的移动能力、可靠的环境感知能力以及能够根据任务目标实时调整行动策略的高层决策能力。在这一系列能力之中，导航是连接“感知”与“执行”的核心环节，它直接决定机器人能否将高层任务意图转化为一条安全、连续、有效的行动过程。

传统移动机器人导航系统通常采用“感知建图、定位估计、路径规划、运动控制”的经典流水线结构\cite{thrun2005probabilistic}。这类方法以 SLAM 类算法\cite{orbslam3_2021,cartographer2016}为核心获取几何地图与位姿估计，其在环境静态、地图充分、传感器标定精确的场景中表现良好，并且具有结构清晰、模块边界明确、理论可解释性强等优点。然而，当机器人面对的是动态变化的真实环境、目标表达不再以精确坐标给出、或者场景中缺少完善几何地图时，传统方法往往会暴露出较强的地图依赖、建模成本高、跨场景泛化能力有限等问题\cite{anderson2018evaluation}。对于需要快速部署、低成本扩展、弱地图依赖的现实任务而言，仅依赖几何地图和人工规则的导航方式已难以完全满足需求。

在此背景下，视觉导航逐渐成为机器人自主移动研究中的重要方向。相较于激光雷达导航或严格依赖环境几何重建的方法，视觉导航直接以图像作为环境主要观测形式，通过学习方法从大量轨迹和图像序列中提取可迁移的导航规律。视觉目标不仅可以用坐标表示，也可以通过目标图像、拓扑节点图像、语义类别，甚至自然语言描述的形式给出。这使得导航任务不再局限于“到达某个几何点”，而更接近于人类在真实环境中的行动方式，即根据当前所见与目标所见之间的关系，不断调整自身朝向和局部路径，逐步接近目标区域。

在学习型视觉导航方法的发展过程中，General Navigation Models（GNM）、Visual Navigation Transformer（ViNT）以及 NoMaD 代表了从判别式局部规划到生成式行为建模的重要演进\cite{gnm2022,vint2023,nomad2023}。GNM 和 ViNT 主要采用视觉编码与时序建模结构，从当前观测和目标观测中直接回归局部路径点或距离估计；NoMaD 则进一步引入动作扩散策略\cite{diffusionpolicy2023}，将机器人未来的局部轨迹视为一个条件生成问题进行建模\cite{nomad2023}。在这一框架下，机器人不再只给出单一步骤上的确定性控制，而是学习“在当前观测和目标条件下，未来一段时间内可能采取的局部运动轨迹分布”。这种生成式策略特别适合解决导航任务中的多模态问题，即在同一观测下，往往存在多条都可行的未来路径，模型需要学会在不确定环境中生成平滑、可达、对目标有牵引作用的轨迹序列。

从研究方法视角看，这种演化还带来一个重要变化：导航模型的价值不再只由“单次预测是否准确”决定，而是越来越取决于其是否能够在闭环系统中长期稳定工作。对于纯离线评测而言，一条局部轨迹只要和标签足够接近就可能被视为有效；但对于真实机器人系统而言，轨迹的推理速度、平滑性、可执行性、对目标的牵引强度以及与控制器的接口兼容性同样重要。换言之，导航研究已经从“只比较模型指标”逐渐转向“同时比较模型、系统与平台三者的耦合质量”。本文的研究正是在这一转向背景下展开的。

然而，现有 NoMaD 及其相关验证工作主要面向轮式平台展开。轮式平台的优势在于动力学模型相对简单、机身姿态变化较小、视觉观测稳定，因而高层导航策略更容易与底层速度控制直接对接。但真实复杂场景中，单纯依靠轮式底盘并不能覆盖全部任务需求。楼梯、门槛、非平整地面、坡面、狭小通道和轻度非结构化区域都对平台通过能力提出了更高要求。与轮式机器人相比，足式机器人，尤其是四足机器人，具有更强的地形适应性与环境穿越能力，因此在复杂环境巡检、安防侦查、灾害搜救和室内外混合场景中具有更高的应用潜力。

足式平台带来能力提升的同时，也显著增加了视觉导航部署难度。首先，四足机器人在行走过程中存在周期性机身起伏、足端冲击与相机抖动，视觉输入的稳定性远低于轮式平台。其次，足式平台通常具有更高的底层控制频率、更严格的动作平滑性要求和更强的运动安全约束，高层导航策略输出的局部轨迹如果缺乏连续性、实时性或可执行性，容易在闭环中放大误差。再次，足式系统往往需要在高层规划、中层控制映射和底层运动执行之间建立明确接口，高层策略不能简单等同于底层关节控制。因此，将已在轮式平台验证过的视觉导航方法迁移到四足平台，并不是“更换底盘”这样简单的工程问题，而是一个涉及视觉表征、轨迹生成、控制解耦和系统闭环验证的综合研究问题。

尤其对于本文关注的 Lite3 平台而言，问题还具有明显的工程约束属性。机器人既要在有限计算资源下完成高层推理，又要满足速度限幅、超时保护、状态机切换和平台通讯等条件；系统既要支持导航模式，也要支持探索模式、多目标任务和仿真/真机后端切换。因此，本文所面对的不是单一算法准确率问题，而是一个典型的“算法进入系统后是否仍然成立”的问题。只有同时处理好模型表达能力与系统接口约束，基于扩散策略的高层视觉导航才有可能真正落地。

\begin{figure}[htbp]
    \centering
    \fbox{\rule{0pt}{55mm}\rule{0.86\linewidth}{0pt}}
    \caption{图位占位：本文研究背景与问题场景示意图（待补）}
    \label{fig:intro-background-placeholder}
\end{figure}

基于上述背景，本文围绕 visualnav-transformer 项目中的 NoMaD 技术链路展开研究，重点关注“基于动作扩散策略的足式机器人视觉目标导航与探索”这一主题。本文的研究工作并不止于复现原始算法，而是进一步面向部署需求和足式场景需求，对 NoMaD 的推理效率、目标引导能力、视觉编码表征能力以及系统集成方式进行系统分析和改进。通过对开题资料、中期研究材料、现有代码实现、仿真验证脚本和真机部署方案的整体梳理，本文试图回答以下几个核心问题。

\begin{enumerate}
    \item 如何从项目代码层面准确理解 NoMaD 中视觉编码器、距离预测头和扩散策略之间的关系，并形成可复现、可训练、可推理的研究基线；
    \item 如何通过 DDIM 采样加速、Classifier-Free Guidance（CFG）引导增强和视觉编码器替换等方法提升 NoMaD 的部署潜力；
    \item 如何将原本主要面向轮式平台的高层视觉导航策略迁移到云深处 Lite3 四足平台上，并构建覆盖离线实验、MuJoCo 仿真、状态机控制和真机部署的完整技术链路。
\end{enumerate}

从研究意义上看，本文的工作一方面有助于完善动作扩散策略在目标导航任务中的工程化理解，为基于生成模型的机器人决策研究提供更接近部署场景的验证路径；另一方面也有助于回答“高层视觉导航策略是否能够跨平台迁移到足式系统”这一具有现实价值的问题。特别是在当前具身智能研究强调“从模型能力走向系统能力”的趋势下，本文不再只关注单一模型离线指标的提升，而更关注算法、系统和平台三者是否能够形成真正闭环、可运行、可分析的整体方案。

进一步地，本文研究还具有明确的方法论意义。围绕 NoMaD 的整理与改进并不是孤立工作，而是为后续类似视觉条件生成策略的部署提供了一种可参考范式：先明确模型内部条件结构，再找出影响部署的主要瓶颈，随后通过统一推理模块和统一实验脚本建立可重复验证链路，最后再进入仿真与真机闭环阶段。这样的研究路径有助于避免“只见模型不见系统”的常见问题，也使论文中的结论更具说服力。

\section{国内外研究现状}

\subsection{传统导航与学习型导航的发展脉络}

从整体发展脉络看，机器人导航研究大致经历了两个主要阶段。第一阶段是以几何建模和概率估计为核心的经典导航范式，其代表性问题包括同步定位与建图（SLAM）、全局路径规划、局部避障和轨迹跟踪。该类方法以环境几何结构为基础，通过激光雷达、深度相机或组合导航传感器估计机器人位姿，再结合栅格地图、拓扑地图或语义地图进行规划与控制。经典方法的优势在于模块边界清晰、误差来源可分析、在规则场景中可获得较高精度和可靠性。

第二阶段则是深度学习推动下的学习型导航阶段。该阶段的核心思想在于：不再完全依赖人工构造的几何中间变量，而是直接从图像观测、历史轨迹和任务目标中学习策略表示\cite{zhu2017targetdriven}。学习型导航通过大规模数据驱动的方式获得跨环境泛化能力，尤其在目标表达灵活性、感知到决策的紧耦合建模以及跨场景迁移方面表现出明显优势。随着视觉表征学习和 Transformer 架构\cite{transformer2017,vit2020}的发展，导航系统逐渐从“局部规则驱动”走向“统一策略模型驱动”，也为本文所关注的通用导航模型奠定了基础。

需要指出的是，学习型导航并不是对经典导航的完全替代。对于大尺度场景、强几何约束场景或高精定位任务，几何方法仍然具有不可忽视的价值。当前国内外研究的一个重要趋势，是将几何先验与学习策略有机结合，使模型既具备学习方法的泛化能力，又能在关键环节保留结构化约束。本文的工作同样遵循这一思路：在高层采用 NoMaD 进行视觉条件下的局部轨迹生成，而在中层控制与平台执行端保留清晰的结构化接口。

\subsection{视觉导航与目标导航研究现状}

视觉导航的核心问题可以概括为：如何根据当前视觉观测与任务目标之间的关系，生成一条能够驱动机器人逐步接近目标的行动策略。早期视觉导航方法多聚焦于 point-goal navigation\cite{anderson2018evaluation}，即目标以相对方向或相对坐标的方式给出，模型需要从图像中恢复足够的环境信息，以支持面向目标点的局部控制。随后，image-goal navigation、object-goal navigation 和 language-goal navigation 等更具表达灵活性的任务形式逐渐受到关注\cite{zhu2017targetdriven,shah2021ving}。此类任务不再要求目标被编码为精确几何位置，而是允许使用目标图像、目标类别或文本描述来指定任务。

图像目标导航尤其适合真实部署场景。对于机器人而言，“到达与某张图像对应的目标位置附近”往往比“到达绝对坐标系中的某个点”更容易与人类任务描述对齐。围绕这一问题，研究者提出了多种目标建模方式，包括直接拼接当前图像与目标图像、构建观测与目标的联合表征、建立视觉记忆库、或者利用多步历史观测推断与目标之间的相对关系。这类方法的关键难点在于：机器人需要理解当前视角与目标视角之间的差异，同时在局部遮挡、光照变化、动态干扰和视角偏移等情况下保持策略稳定。

在视觉目标导航进一步发展过程中，拓扑导航与视觉记忆方法逐渐成为重要支撑\cite{savinov2018sptm,chaplot2020nts}。其基本思想不是恢复完整精确地图，而是保存若干具有代表性的节点图像，并利用节点之间的顺序关系、可达关系或语义联系来进行导航\cite{shah2021ving}。与精确几何地图相比，拓扑表示更轻量，也更适合与图像目标表达直接结合。对于本文研究的系统而言，topomap 节点既是部署时的候选目标集合，也是高层 NoMaD 进行局部子目标选择的重要中间载体。这种方式兼顾了视觉目标导航的灵活性和系统部署的可操作性。

\subsection{通用导航模型与生成式导航研究现状}

围绕“能否用统一模型处理多种导航场景”这一问题，通用导航模型近年来取得了快速发展。GNM 强调通过统一的图像观测和目标表征学习可迁移的局部导航能力\cite{gnm2022}；ViNT 则进一步利用 Transformer 结构\cite{transformer2017}对观测序列和目标图像进行联合建模，使模型能够在更长时序范围内理解任务条件与历史视觉上下文之间的关系\cite{vint2023}。上述工作的重要共同点在于，它们都弱化了具体底盘控制命令在高层策略中的主导作用，转而强调输出更抽象的局部路径点或轨迹表示。

这一思路为跨平台迁移提供了重要启发。如果高层策略直接输出轮速、关节命令等与平台强耦合的控制量，那么模型迁移到新平台时往往需要彻底重新设计；而如果高层策略输出局部 waypoint、局部子目标或轨迹片段，那么中层控制器可以根据平台特性完成再映射。本文后续能够将 NoMaD 从原始轮式场景扩展到 Lite3 四足平台，本质上也是依赖于这一“高层输出轨迹、中层负责执行”的分层思想。

进一步地，随着生成模型的发展，导航研究逐渐从“判别式输出一个最优动作”转向“生成式建模一组可能未来轨迹”。生成式导航方法特别适合处理局部导航中的多解问题，例如当机器人面对拐角、门洞或开阔区域时，往往存在多条都可行的局部路径。如果高层模型只采用单一回归输出，很容易被训练数据中的平均效应拖向不够清晰的中间解；而生成式方法则有潜力更自然地表示多模态未来行为分布。NoMaD 将这一思想系统引入目标导航，形成了以扩散策略为核心的视觉条件轨迹生成方案\cite{nomad2023}。

\subsection{扩散模型与扩散策略研究现状}

扩散模型最初主要在图像生成领域获得广泛影响，其代表性基础工作包括 DDPM 和 DDIM\cite{ddpm2020,ddim2020}。这类方法通过前向加噪和反向去噪过程学习复杂样本分布，与传统自回归模型或判别式回归模型相比，在建模多模态分布和生成平滑样本方面表现出良好能力。由于机器人策略学习本身同样面临多模态输出问题，研究者开始将扩散模型引入轨迹规划、离线决策、机械臂操控与视觉运动控制等具身智能任务中，逐步形成所谓扩散策略\cite{diffuser2022,diffusionpolicy2023}。

扩散策略的关键优势在于，它不是直接对某一时刻的动作进行点估计，而是在条件信息的约束下生成一整段动作序列或轨迹序列。若从具身智能领域的演化脉络来看，Diffuser 较早以“轨迹去噪”的形式把扩散思想系统引入机器人决策与规划问题，Diffusion Policy 则进一步将动作扩散明确用于视觉运动策略学习，并证明这种方式可以在多任务操控中稳定学习时序动作分布\cite{diffuser2022,diffusionpolicy2023}。对于机器人导航而言，这种建模方式有三点重要意义。第一，它更适合表示局部导航中的多种合理轨迹；第二，它天然能够输出具有时序结构的轨迹片段，从而便于与中层控制器对接；第三，它为部署阶段引入快速采样、条件引导等推理策略提供了明确接口。NoMaD 正是利用扩散策略生成未来一段时间内的 waypoint 序列，再配合距离预测头与 topomap 候选节点选择，实现目标导航与探索模式的统一建模\cite{nomad2023}。

不过，扩散策略也带来了新的部署挑战。标准 DDPM 采样通常需要较多迭代步数，容易在机器人在线推理中产生较大延迟；同时，模型在有条件和无条件两种模式下的输出差异，也意味着如何增强目标条件的牵引能力成为一个值得研究的问题。针对这些问题，DDIM 提供了更快的近似采样方式，CFG 则提供了增强条件约束的简洁机制\cite{ddim2020,cfg2022}。围绕这两者展开部署导向实验，是本文算法部分的重要内容之一。

\subsection{足式机器人导航与部署研究现状}

与轮式导航相比，足式机器人导航研究更长期聚焦于“如何稳健运动”而非“如何高层导航”。许多足式研究工作把主要精力放在步态规划、平衡控制、地形估计与本体控制上\cite{hwangbo2019anymal,lee2020learning,rl_locomotion_rudin2022}，而将高层导航作为相对独立、甚至较为简化的模块处理。随着四足平台逐步进入巡检、安防和复杂环境作业场景，高层导航的角色正在变得越来越重要。机器人不仅要能走、能稳，还要能够根据视觉目标进行自主决策、在未知区域中探索、在局部失败后恢复并继续任务。

当前足式平台上的视觉导航研究仍面临几个典型困难。首先，感知端存在显著扰动，机体姿态变化导致相机画面频繁偏移和抖动；其次，决策端存在更强的实时性约束，高层轨迹如果生成过慢，将直接影响闭环控制稳定性；再次，执行端存在更复杂的接口映射，高层轨迹必须经过中层控制器转化为可执行的参考量，才能交给底层运动系统安全执行。这意味着，在足式平台上验证视觉导航方法时，仅比较离线准确率或单步前向推进指标是不够的，还必须从系统角度考察策略输出是否连续、是否稳定、是否真正可落地。

因此，国内外足式导航研究的一个现实问题在于：虽然已经出现一些将学习型感知与高层决策结合到足式系统中的探索，但从统一算法框架、分层控制设计、仿真验证到真机部署的完整闭环链条仍然相对稀缺。本文正是在这一空缺位置上展开工作，希望构建一条从 NoMaD 算法复现、算法优化，到 Lite3 仿真与真机部署的完整技术路径。

\begin{figure}[htbp]
    \centering
    \fbox{\rule{0pt}{55mm}\rule{0.86\linewidth}{0pt}}
    \caption{图位占位：视觉导航相关工作演化脉络图（待补）}
    \label{fig:intro-related-work-placeholder}
\end{figure}

\subsection{现有研究不足与本文切入点}

综合以上分析，可以看出当前研究仍存在四个层面的不足。第一，视觉目标导航领域虽然已经形成了较成熟的学习型方法谱系，但面向足式系统的生成式局部轨迹策略研究仍然不足，特别是缺少围绕闭环部署需求展开的系统性分析。第二，NoMaD 虽然已经证明了扩散策略在目标导航中的潜力，但其原始验证主要集中于轮式平台，尚未系统回答高层扩散策略是否能稳定服务于四足机器人。第三，关于推理加速、目标引导强度和视觉编码器替换等问题，现有工作往往以单点实验方式呈现，缺少在统一项目框架下可复现、可比较的系统实验。第四，从算法到系统再到平台的衔接仍然是薄弱环节，很多论文在离线指标层面效果良好，但没有进一步回答“如何真正接入仿真和真机闭环”。

针对这些不足，本文的切入点可以概括为以下三点。其一，从项目现有代码出发，对 NoMaD 的训练结构、推理结构与部署接口进行系统整理，明确视觉编码器、距离预测头、扩散策略和 goal mask 机制的内部关系。其二，围绕部署可行性对 NoMaD 进行定向优化，重点考察 DDIM、CFG 和视觉编码器升级对推理效率、轨迹质量和目标牵引能力的影响。其三，面向云深处 Lite3 四足平台构建分层闭环系统，将高层 NoMaD 推理、中层状态机与控制映射、底层运动接口有机结合，并通过 MuJoCo 仿真与真机部署验证形成完整研究链条。

由此可见，本文关注的并不仅是“某个指标能否提升”，而是“基于动作扩散策略的视觉目标导航，能否在足式机器人上形成一套可复现、可分析、可优化、可部署的系统方案”。这一问题既具有理论研究意义，也具有明确的工程应用价值。

\subsection{本文研究问题的分层定义}

为了避免把不同层次的问题混写在一起，本文将研究对象明确划分为算法层、系统层和平台层三个层级。算法层关注的是高层视觉导航策略本身，即在给定历史视觉观测和目标条件的情况下，模型能否生成方向合理、平滑稳定、与目标具有明确关联的局部轨迹；系统层关注的是这些高层轨迹如何被统一推理模块、状态机、导航主机和中层控制接口接纳，从而构成一个能够连续运行的闭环链路；平台层关注的则是该闭环链路在 Lite3 四足机器人与其配套计算平台上是否能够满足实时性、安全性和任务组织需求。

这三个层级之间既有明确分工，又存在强耦合关系。算法层若只追求离线统计指标提升，而忽视输出轨迹的平滑性与执行可行性，系统层就会在状态切换或中层控制映射阶段暴露问题；系统层若缺少统一推理接口和任务调度机制，则很难稳定复用算法层的研究成果；平台层若存在算力瓶颈、通信时延或安全约束，则会反过来要求算法与系统主动压缩推理延迟、限制输出幅度并引入恢复机制。换言之，本文并不是单纯研究一个更“准”的模型，而是在研究一个高层视觉导航策略如何真正进入机器人系统并保持有效。

基于这样的分层定义，本文后续章节的组织逻辑也更加清晰。第二章讨论的是理解问题所需的共通知识，包括导航任务形式、扩散建模原理、轮足平台差异和 Lite3 的工程约束；第三章讨论的是算法层问题，即 NoMaD 的结构、训练与离线优化实验；第四章进一步转向系统层，研究推理模块、状态机、导航主机和 MuJoCo 闭环验证；第五章再进入平台层，关注 Orin 上的真机部署、任务运行与可视化记录。通过这种分层写法，本文希望避免“离线指标提升就等于系统成功”这一常见误判，使论文结论与实际研究边界保持一致。

\subsection{本文研究的技术路线与验证思路}

围绕上述分层问题，本文采用了一条由“项目梳理”逐步走向“系统闭环”的技术路线。首先，在研究起点上并不预设某个全新算法，而是立足于现有 \textit{visualnav-transformer} 项目，对其中与 NoMaD 相关的配置、训练脚本、推理入口、实验结果和部署代码进行统一梳理。这样做的目的，是确保论文中的每一项结论都可以回溯到真实项目实现，而不是停留在脱离代码基础的概念分析上。对于毕业设计而言，这一步尤为关键，因为只有先建立从代码到论文叙述的一一对应关系，后续的优化与扩展才具有可验证性。

在此基础上，本文把后续研究工作进一步划分为“离线筛选”和“闭环验证”两个阶段。离线筛选阶段主要面向算法层，重点分析 NoMaD 的内部结构、goal mask 机制、距离预测头与扩散策略头的耦合关系，并围绕 DDIM、CFG 和视觉编码器升级设计可统计、可比较的实验。该阶段的任务不是直接给出最终部署方案，而是通过统一实验框架回答哪些因素更可能成为部署瓶颈、哪些参数更值得进入下一阶段验证。闭环验证阶段则进一步转向系统层与平台层，将算法层得到的候选配置真正放入 MuJoCo 仿真和 Lite3 真机系统中，检验其在状态机、控制映射、通信接口和算力约束下是否仍能保持有效。

这种技术路线的一个重要特点，是始终强调“证据链完整性”。换言之，本文并不把某一次训练结果或某一组离线指标孤立地看作结论，而是要求结论能够沿着“模型结构解释一组指标变化，指标变化支持候选方案筛选，候选方案再进入闭环系统验证”这一链条逐步成立。这样一来，论文既能够保留算法研究的深度，也能够体现毕业设计应有的系统实现完整性。对最终成果来说，真正有价值的并不只是单个数字的提高，而是算法优化、系统组织和平台部署三者之间形成了彼此支撑的研究闭环。

\section{本文结构安排}

围绕上述研究问题，本文以 visualnav-transformer 项目为基础，以 NoMaD 为高层核心算法，以 Lite3 四足平台为主要部署对象，按照“预备知识、算法分析、系统设计、仿真验证、真机实验”的主线展开全文。本文的主要研究内容与工作贡献可归纳如下。

\begin{enumerate}
    \item 对 NoMaD 相关代码、配置、训练链路和推理链路进行系统梳理，完成从项目实现到论文表述之间的结构化映射，明确高层视觉编码器、距离预测分支与扩散策略分支之间的关系。
    \item 面向部署需求对原始 NoMaD 进行算法级优化，重点包括基于 DDIM 的采样加速、基于 CFG 的目标条件增强，以及基于 backbone-suite 的视觉编码器统一替换与公平对照实验设计。
    \item 面向足式平台构建由推理模块、状态机、导航主机、中层控制接口和底层执行接口组成的分层闭环系统，使高层视觉导航策略能够在 MuJoCo 仿真和真机系统中统一运行。
    \item 以云深处 Lite3 四足平台为载体，形成包含视觉数据采集、topomap 构建、任务切换、运行日志保存、图像可视化与调试接口在内的完整部署流程，并分析算法与系统在真实平台上的适用性与局限性。
\end{enumerate}

为了更直观地说明本文工作的组织方式，还可以将其技术路线概括为四个连续阶段。第一阶段是项目理解与基线建立阶段，主要工作是阅读并整理 visualnav-transformer 项目中的 NoMaD 相关实现、配置文件、训练脚本和部署脚本，明确后续实验的真实起点。第二阶段是算法优化与离线实验阶段，围绕推理加速、目标引导和视觉编码器替换完成统一实验设计。第三阶段是系统设计与仿真闭环阶段，构建推理模块、状态机与导航主机，使高层策略能够在 MuJoCo 中完成完整任务流程。第四阶段是真机部署阶段，在 Orin 与 Lite3 平台上验证导航模式、探索模式、图像可视化与运行记录机制。

从论文撰写角度看，这样的技术路线也决定了全文不能只按照“先讲算法、再给几个结果”的传统方式组织，而必须同时体现出项目整理、方法优化、系统设计和平台验证之间的因果关系。也正因此，本文在每一章中都尽量把“为什么这样设计”“这一阶段回答了什么问题”“它与下一阶段是什么关系”交代清楚，以保持全文逻辑连贯。

在章节安排上，全文共分为六章，各章内容如下。

\begin{enumerate}
    \item 第一章为绪论，介绍课题研究背景、国内外研究现状、本文要解决的问题以及整体研究思路。
    \item 第二章为预备知识，重点说明导航任务、视觉导航、目标导航、拓扑导航、扩散模型、扩散策略、轮式与足式平台差异，以及云深处 Lite3 平台的相关参数与工程约束。
    \item 第三章为算法设计与实验，围绕 NoMaD 的算法复现、输入输出形式、训练机制、推理机制以及 DDIM、CFG、视觉编码器升级等优化实验展开分析，并给出离线统计实验结果。
    \item 第四章为算法框架设计与仿真环境验证，介绍推理模块、状态机、导航主机和分层闭环系统的设计方法，并在 MuJoCo 中对 Lite3 导航与探索流程进行验证。
    \item 第五章为真机实验，面向 Orin 与 Lite3 平台给出部署流程、任务运行方式、可视化与图像保存方案，并分析真机运行中的效果与问题。
    \item 第六章为总结与展望，对全文工作进行概括，分析当前方法的不足，并讨论后续在多目标导航、视觉记忆增强、跨平台扩展和真机长期运行等方向上的改进空间。
\end{enumerate}

从全文结构关系上看，这六个章节并不是简单并列的内容拼接，而是围绕同一研究主线逐层递进的。第一章负责提出问题、明确研究边界和总体路线；第二章负责建立理解这些问题所需的共同语言；第三章则在这一基础上回答“高层算法如何设计与优化”；第四章和第五章继续回答“这些高层算法如何真正进入仿真和真机系统，并在平台约束下保持有效”；第六章则对前述结果进行统一回收和再判断。这样的组织方式使得本文既保持了算法研究的理论深度，也能够体现毕业设计中系统实现与工程落地的完整性。

进一步地，本文各章输出还可以被理解为三类不同性质的研究成果。第二、三章主要产出的是模型理解与算法证据，包括对 NoMaD 结构的解析、对 DDIM 和 CFG 的部署导向分析、以及对视觉编码器升级路线的阶段性结论；第四章主要产出的是系统证据，即统一推理模块、状态机、导航主机和 MuJoCo 闭环链条能否把第三章的候选方案稳定接入；第五章主要产出的是平台证据，即在 Lite3 真机与 Orin 边缘计算平台上，这套系统是否具备真正运行和记录的能力。将这些输出层次区分开来，有助于避免在论文中把“算法上有潜力”“系统上可运行”和“平台上可部署”混写为同一种结论。

这种区分不仅有助于论文结构清晰，也直接影响结论表述的严谨性。举例来说，一个视觉编码器在离线评测中表现更好，只能说明它在当前统一训练与评测框架中更具潜力；只有当它进一步经过系统接入和平台验证之后，才有资格被写成更适合部署的方案。类似地，一个推理加速策略若仅在延迟指标上占优，但闭环后导致状态切换不稳定，也不能被视作整体更优。正因为本文始终坚持这种分层判断原则，全文才会采用“先离线筛选、再闭环验证、最后真机确认”的组织方式。

从毕业论文写作要求看，这样的组织方式还有一个额外优点，即能够把“完成了哪些工作”和“这些工作分别证明了什么”清晰地区分开来。前者对应工作量与实现完整性，后者对应学术结论与工程结论的边界。二者结合，才能使全文既有足够扎实的内容厚度，也有经得起追问的论证结构。

这也意味着，本文在后续章节中会始终坚持“问题提出、方法设计、证据验证、边界说明”四步展开，而不是只罗列实现过程或只堆叠实验数字。这一写作原则同样对应了本文从开题、中期推进到最终系统落地的真实工作推进顺序，也是全文能够形成完整研究闭环、使章节之间的衔接更具可解释性并便于答辩展开的重要原因。

除了章节安排之外，本文的研究输出还可以从“算法结果”“系统结果”和“平台结果”三个层面来理解。算法结果关注的是 NoMaD 高层策略在离线条件下的可解释优化，包括 DDIM、CFG 和视觉编码器替换是否真的改变了推理速度和轨迹生成质量；系统结果关注的是这些算法优化是否能够被统一推理模块、状态机和导航主机稳定接纳，形成可复用的运行链路；平台结果则关注 MuJoCo 与 Lite3 真机是否能够在相同高层逻辑下完成任务。将这三层输出区分开来，有助于避免把某一层的局部改进误写成对整套系统的全面结论。

也正因如此，本文在表述上会尽量避免“某个方法最好”这类脱离验证层次的绝对化结论，而更强调“某种方法在哪个阶段显示出更高潜力”“某种配置在什么条件下更适合进入下一阶段验证”。这种写法既更符合工程研究的真实推进过程，也更便于后续在答辩和论文修改中持续补充证据。

总体而言，本文并不是将已有 NoMaD 工作简单迁移到另一台机器人上，而是试图完成一次从“算法论文结果”到“可复现系统方案”的系统化整理与扩展。通过这一研究过程，本文希望为生成式导航策略在足式机器人上的应用提供更加扎实的实验依据和工程实现参考。

\begin{figure}[htbp]
    \centering
    \fbox{\rule{0pt}{55mm}\rule{0.86\linewidth}{0pt}}
    \caption{图位占位：本文总体技术路线与章节关系图（待补）}
    \label{fig:intro-roadmap-placeholder}
\end{figure}
